\documentclass[10pt]{article}
\usepackage[T1]{fontenc}
\usepackage[utf8]{inputenc}
\usepackage[french]{babel}
\usepackage{amsmath,amssymb,mathrsfs,mathpazo}
\usepackage{geometry}
\geometry{a4paper,margin=2cm}
\usepackage{graphicx,array,tabularx,tabularray,colortbl,multirow,multicol,booktabs,enumitem,pifont,ulem,fancybox,diagbox,cancel}
\usepackage{pstricks,pstricks-add,pst-plot,pst-tree,pst-node,pst-text,pst-func,pst-math,pst-bspline,pst-3dplot,pst-eucl,pst-solides3d}
\let\pscalculate\undefined
\usepackage{pst-bezier}
\usepackage{tikz,pgfplots,tkz-tab}
\usetikzlibrary{arrows,arrows.meta,patterns,calc,positioning,trees,decorations.markings,shapes.geometric,shapes.misc,shapes.arrows,intersections,decorations.pathmorphing,decorations.pathreplacing,backgrounds,shadows,babel}
\pgfplotsset{compat=1.18}
\usepgfplotslibrary{fillbetween}
\usepackage[np]{numprint}
\usepackage{listings,makecell,pas-tableur,twoopt}
\let\barre\undefined
\usepackage[pstricks]{tablvar}
\definecolor{darkgreen}{RGB}{0,100,0}
\providecommand{\red}{\color{red}}
\providecommand{\blue}{\color{blue}}
\providecommand{\green}{\color{green}}
\providecommand{\black}{\color{black}}
\providecommand{\white}{\color{white}}
\providecommand{\gray}{\color{gray}}
\providecommand{\cyan}{\color{cyan}}
\providecommand{\magenta}{\color{magenta}}
\providecommand{\yellow}{\color{yellow}}
\providecommand{\violet}{\color{violet}}
\providecommand{\orange}{\color{orange}}
\providecommand{\euro}{\texteuro}
\providecommand{\EUR}{\texteuro}
\providecommand{\ieme}{\textsuperscript{e}}
\providecommand{\degres}{\ensuremath{^\circ}}
\providecommand{\degree}{\ensuremath{^\circ}}
\providecommand{\vv}[1]{\overrightarrow{#1}}
\providecommand{\Coord}[1]{#1}
\providecommand{\decosix}{}
\providecommand{\starredbullet}{\ensuremath{\bullet}}
\pagestyle{empty}
\setlength{\parindent}{0pt}
\usepackage{adjustbox}\DeclareUnicodeCharacter{2236}{:}\DeclareUnicodeCharacter{2190}{\ensuremath{\leftarrow}}
\providecommand{\R}{\mathbb{R}}
\providecommand{\N}{\mathbb{N}}
\providecommand{\Z}{\mathbb{Z}}
\providecommand{\Q}{\mathbb{Q}}
\providecommand{\D}{\mathbb{D}}
\providecommand{\C}{\mathbb{C}}
\providecommand{\e}{{\mathrm{e}}}
\providecommand{\vect}[1]{\overrightarrow{#1}}
\providecommand{\vectt}[1]{\overrightarrow{\mathrm{#1}}}
\providecommand{\vvt}[1]{\overrightarrow{\mathrm{#1}}}
\providecommand{\np}[2][]{{#2~#1}}
\providecommand{\vv}[1]{{\overrightarrow{#1}}}
\renewcommand{\d}{{\mathrm{d}}}
\providecommand{\degres}{{^\circ}}
\providecommand{\degre}{{^\circ}}
\providecommand{\blue}{\color{blue}}
\providecommand{\red}{\color{red}}
\providecommand{\enskip}{\quad}
\providecommand{\llbracket}{\lbrack\!\lbrack}
\providecommand{\rrbracket}{\rbrack\!\rbrack}
\providecommand{\euro}{{\text{€}}}
\providecommand{\no}{{\text{n°}}}
\providecommand{\interval}[2]{{[#1;#2]}}
\providecommand{\intervFO}[2]{{[#1;#2[}}
\providecommand{\dotfill}{\ldots}

\providecommand{\ds}{\displaystyle}
\providecommand{\textding}[1]{#1}
\providecommand{\Oij}{\(\mathrm{O};\vec{\imath},\vec{\jmath}\)}
\providecommand{\Oijk}{\(\mathrm{O};\vec{\imath},\vec{\jmath},\vec{k}\)}
\providecommand{\Ouv}{\(\mathrm{O};\vec{u},\vec{v}\)}
\providecommand{\intervFF}[2]{\ensuremath{[#1;#2]}}
\renewcommand{\euro}{\ifmmode\text{\texteuro}\else\texteuro\fi}
\DeclareUnicodeCharacter{00B0}{\ifmmode^{\circ}\else\textdegree\fi}
\DeclareUnicodeCharacter{20AC}{\ifmmode\text{\texteuro}\else\texteuro\fi}
% Énoncés : Éducation nationale. Transcription et corrigés : APMEP, Denis Vergès et contributeurs. Découpage et classement : CoopMaths, Sébastien Lozano, Éric Elter et contributeurs. Adaptation de la présentation : Maths Explorer.
% CC BY-SA : https://coopmaths.fr/www/about/
% https://coopmaths.fr/alea/static/bac/2026/tex/bac_2026_06_sujet1_polynesie_4.tex
\begin{document}
\noindent\textit{APMEP / CoopMaths. CC BY-SA. Présentation adaptée par Maths Explorer.}\par\medskip
L'espace est muni d'un repère orthonormé \Oij.

On considère les points A$(1~;2~;~-1)$, B$(0~;3~;~2)$, C$(-2~;4~;~0)$ et D$(8~;2~;~-11)$.

\begin{enumerate}
\item
\begin{enumerate}[label=\textbf{\alph*.}]
\item Démontrer que les points A, B et C définissent un plan.
\item Démontrer que le vecteur $\vect{n}\begin{pmatrix}5\\8\\-1\end{pmatrix}$ est un vecteur normal au plan (ABC).
\item En déduire qu'une équation cartésienne du plan (ABC) est $5x+ 8y - z - 22 = 0$.
\item Vérifier que le point D n'appartient pas au plan $(ABC)$.
\end{enumerate}

\item
	\begin{enumerate}[label=\textbf{\alph*.}]
		\item On considère la droite $\Delta$, orthogonale au plan (ABC) et passant par D.

Justifier qu'une représentation paramétrique de la droite $\Delta$ est~:
\[
\begin{cases}
x = \phantom{-}8+ 5t\\
y = \phantom{-}2+ 8t\\
z = -11- t
\end{cases}, \quad \text{avec } t \in \R.
\]
		\item On appelle E le point d'intersection de la droite $\Delta$ et du plan (ABC).

Démontrer que le point E a pour coordonnées $\left(\dfrac{11}{2}~;~-2~;~-\dfrac{21}{2}\right)$.
		\item En déduire la valeur exacte de la longueur DE, distance du point D au plan (ABC).
	\end{enumerate}
\item On note H le point de coordonnées $\left(-\dfrac{2}{3}~;\dfrac{10}{3}~;\dfrac{4}{3}\right)$.

On admet qu'une représentation paramétrique de la droite (BC) est~:

\[
\begin{cases}
x = \phantom{3}-2k\\
y = 3+k\\
z = 2-2k
\end{cases}, \quad \text{avec } k \in \R.
\]

Avec un logiciel de calcul formel, on obtient les résultats suivants~:

\begin{center}
\begin{tabular}{|p{8.5cm}|}\hline
\textbf{Entrée 1~:} $(-2/3-1)*(-2)+(10/3-2)*1+(4/3-(-1))*(-2)$\\
\textbf{Sortie 1~:} $\rightarrow 0$\\
\hline
\textbf{Entrée 2~:} RésoudreSystème$(-2k=-2/3,\ 3+k=10/3,\ 2-2k=4/3)$\\
\textbf{Sortie 2~:} $\rightarrow k = 1/3$\\
\hline
\textbf{Entrée 3~:} RacineCarrée$\big((-2/3-1)^2+(10/3-2)^2+(4/3+1)^2\big)$\\
\textbf{Sortie 3~:} $\rightarrow$ RacineCarrée$(10)$\\
\hline
\textbf{Entrée 4~:} RacineCarrée$\big((-2-0)^2+(4-3)^2+(0-2)^2\big)$\\
\textbf{Sortie 4~:} $\rightarrow 3$\\
\hline
\end{tabular}
\end{center}

	\begin{enumerate}[label=\textbf{\alph*.}]
		\item En utilisant les deux premiers résultats obtenus avec la feuille de calculs ci-dessus, justifier que le point H est le projeté orthogonal du point A sur la droite (BC).
		\item En exploitant les autres résultats obtenus avec la feuille de calculs ci-dessus, déterminer le volume du tétraèdre ABCD.

On rappelle que le volume $\mathcal{V}$ d'un tétraèdre est donné par $\mathcal{V} = \dfrac{1}{3}\times\mathcal{B}\times h$, où $\mathcal{B}$ est l'aire d'une base du tétraèdre et $h$ la hauteur relative à cette base.
	\end{enumerate}
\end{enumerate}

\end{document}
